% TOP_PROBLEM: {"problemId": "problem.polynomial-bound-for-limit-cycles-of-planar-polynomial-vector-fields", "canonicalTitle": "Polynomial Bound for Limit Cycles of Planar Polynomial Vector Fields", "releaseRank": 125, "primaryDomain": "probability_ergodic_dynamics", "primaryDomainLabel": "Probability, ergodic theory & dynamics", "displayStatus": "open", "releaseStatus": "open"}
% !TeX program = lualatex
% Definition and sources only; this file is not an LLM solution attempt.
\documentclass[11pt]{article}
\usepackage[margin=25mm]{geometry}
\usepackage{fontspec}
\setmainfont{Latin Modern Roman}
\usepackage{amsmath,amssymb,mathtools}
\usepackage{unicode-math}
\setmathfont{Latin Modern Math}
\usepackage{xurl,hyperref}
\hypersetup{colorlinks=true,urlcolor=blue}
\setcounter{secnumdepth}{0}
\setlength{\parindent}{0pt}
\setlength{\parskip}{5pt}
\setlength{\emergencystretch}{3em}
\begin{document}
\section*{\texorpdfstring{Polynomial Bound for Limit Cycles of Planar Polynomial Vector Fields}{Polynomial Bound for Limit Cycles of Planar Polynomial Vector Fields}}\label{125}
\subsection{Definitions and mathematical statement}
For real polynomials $P,Q$ of degree at most $d$, a limit cycle of $\dot x=P(x,y)$, $\dot y=Q(x,y)$ is an isolated periodic orbit. The target asks whether there is one exponent $q>0$ with
\[
 \#\{\text{limit cycles of }(P,Q)\}\le d^q
\]
for every degree $d$ and every such system, with the nontrivial degree range $d\ge2$. The exponent is independent of all coefficients and of $d$. A finite but faster-than-polynomial degree bound would not answer this stronger quantitative question.
\par\smallskip\textit{Formulation and concept references:} \href{https://www.fim.uni-passau.de/fileadmin/dokumente/fakultaeten/fim/lehrstuhl/muller/SmaleProblems1998.pdf}{[S1]}.

\subsection{Short English statement}
Can the number of limit cycles of every planar polynomial system be bounded by one universal power of its degree?
\subsection{Sources}
\begin{itemize}
\item[S1]Steve Smale, Mathematical Problems for the Next Century, 1998.
\url{https://www.fim.uni-passau.de/fileadmin/dokumente/fakultaeten/fim/lehrstuhl/muller/SmaleProblems1998.pdf}.
\textit{Formulation source.}
\end{itemize}
\end{document}
